%======================================================================
\section{Track MFG --- Manufacturing outputs}\label{sec:mfg}
%======================================================================

BruceWorx parts go to three processes: FDM printing (PETG/ABS), CNC machining and sheet metal / laser.
FreeCAD already exports 3MF (Mesh workbench) and DXF (Draft importers), and its CAM workbench ships
post-processors, including GRBL and LinuxCNC. This track connects those to the model and to the release.

\begin{epic}{MFG-01}{3D-print prep}{M3}{\stTodo}{(new)}
\begin{itemize}
  \item A \textbf{print orientation} property on each part (build direction plus a rotation about it). It is
        shown as a build-plate gizmo and used by the overhang map (\eref{MOD-11}) and the orthotropic
        material (\eref{SIM-03}).
  \item \textbf{3MF export} with orientation applied, units, part name \code{P/N-REV}, material and colour
        metadata.
  \item \textbf{Slicer handoff}: open the 3MF directly in OrcaSlicer, Bambu Studio or PrusaSlicer, configurable,
        with a per-material profile hint.
  \item \textbf{Insert and print-hole presets}: heat-set insert holes by insert size, plus clearance and
        horizontal-hole teardrop shapes (from \eref{MOD-05}).
\end{itemize}
\emph{Acceptance:} a part opens in the slicer already oriented, with its name and material set.
\end{epic}

\paragraph{Theory: choosing an orientation.} For a candidate build direction $\mathbf b$, score
\[
  J(\mathbf b)=w_1\!\!\sum_{f:\ \mathbf n_f\cdot(-\mathbf b)>\cos\theta_c}\!\! A_f
  \;+\;w_2\,h(\mathbf b)\;+\;w_3\,\big|\boldsymbol\sigma_{\text{dominant}}\cdot\mathbf b\big| ,
\]
which adds supported area, build height and the alignment of the main load with the weak Z axis (from a
static study, if one exists). Evaluating the faces of the convex hull plus the principal axes as candidates,
and taking the minimum, gives a suggested orientation that the user can accept or override.

\begin{epic}{MFG-02}{CNC CAM}{M3}{\stTodo}{(new)}
Use FreeCAD CAM with BruceWorx defaults: a \textbf{setup template} (stock from the bounding box plus
allowance, work offset at a datum corner), a \textbf{tool library} matching the shop's tools, operation
templates (face, adaptive clearing, profile, drill and chamfer), and a \textbf{post-processor} for the target
machine (open question).
\emph{Acceptance:} a corpus part posts G-code that simulates clean in the CAM simulator.
\end{epic}

\begin{epic}{MFG-03}{Sheet-metal flat DXF}{M3}{\stTodo}{(new)}
Export each sheet-metal part's flat pattern (\eref{MOD-06}) as DXF, with layers \code{CUT} (outline and
holes), \code{BEND\_UP}/\code{BEND\_DOWN} (bend lines) and \code{ETCH} (part number), and a bend note
block. File name \code{P/N-REV.dxf}. Units are mm, and splines are converted to arcs where the shop's CAM
requires it.
\emph{Acceptance:} the laser shop accepts the DXF without edits.
\end{epic}

\begin{epic}{MFG-04}{Release package}{M2}{\stTodo}{(new)}
When a revision is released (\eref{PDM-05}), build a package: STEP AP242 per part and assembly, the PDF
drawing, DXF flats, 3MF prints and a CSV BOM, plus a manifest with SHA-256 checksums. The package is attached
to the revision in the vault.
\emph{Acceptance:} rebuilding the package from the released files gives identical checksums for every
deterministic output.
\end{epic}

\paragraph{Note on determinism.} STEP and PDF writers embed timestamps. The package normalises them, for
example by setting the STEP header time to the release time, so the checksum identifies \emph{content}.
A vendor who receives \code{LE-02001-A} can then verify that the files are exactly the released ones.
